The TunSociety project resulted in a working community platform adapted to the needs of Tunisian student clubs, associations, and small organized groups. Its purpose is deliberately narrower than a public social network. The application gives members a place to communicate and build connections, while giving responsible actors the tools required to supervise behavior and keep decisions traceable.

The analysis phase clarified the main actors and the value of each workflow. Members need everyday features such as profiles, posts, comments, reactions, member search, friend requests, direct messages, and notifications. Moderators and administrators need a different set of capabilities: content review, user supervision, role management, warnings, freezes, appeals, risk summaries, and audit logs. Designing the system around these two families of needs helped keep the product focused on controlled community life.

Technically, the project was implemented as a full-stack web application. Angular handles the client interface, ASP.NET Core exposes the backend API, Entity Framework Core maps the domain model to MySQL, and JWT tokens protect private routes. The backend remains the trusted layer for authorization, validation, moderation, and administrative actions, while the frontend concentrates on navigation, interaction, and responsive presentation.

The local AI-assisted moderation module is one of the main specific contributions of the project. By using Ollama and Gemma, the system can evaluate submitted text and store a moderation record containing categories, a score, an action, and a reason. This mechanism supports moderators and administrators without replacing human judgment. It also strengthens traceability because decisions can be reviewed later instead of disappearing after the content is processed.

The current version proves the feasibility of combining social interaction and governance inside the same academic project. It also leaves clear directions for improvement: real-time messaging, a progressive web application experience, richer moderation timelines, configurable thresholds, multilingual moderation for Arabic, French, English, and Tunisian dialect, analytics, OAuth integration, deployment automation, and stronger automated tests.

Overall, TunSociety achieved its main objective. It delivers a maintainable web platform where communication, role-based governance, local moderation support, and administrative traceability work together for a clearly identified community context.
